\chapter{Residual-source descent, unbounded exponent tails, and
period-driven cycle exclusion}
{\small\sffamily Source: \nolinkurl{repository/collatz_reconstruction/research_program/residual_splice_20260914/note.md}.\par}


\textbf{14 September 2026 --- continuation in
\nolinkurl{KokunoYumeto/collatz-workbench}.}

Source checkpoint: \nolinkurl{32b8f3b9cc6dcfc4387ff783c4f4a6f6a365f717},
PR \#2. The original stopped-affine, cycle-relative, and integration
files remain unchanged. This note and \nolinkurl{integral_blocks.md}
supply the new arguments. The verification records finite execution, not
independent external mathematical review or a global Collatz proof.

\section{1. Actual source and first-descent cells}

Use the original map
\SourceMath{T(n)=(3n+1)/2^{\nu_2(3n+1)},\qquad n\in X=\{1,3,5,\ldots\}.}{}
For a nonempty exponent word \(p=(a_1,\ldots,a_m)\), keep
\SourceMath{A_j=\sum_{i\le j}a_i,\quad L_j=3^j,\quad U_j=2^{A_j},\quad
 C_j=\sum_{i=1}^j3^{j-i}2^{A_{i-1}},\quad D_j=U_j-L_j.}{} Write
\((A,L,U,C,D)\) for the final data. The predecessor's original legal
cylinder is \SourceMath{n=r+2Ut,\quad t\ge0,\qquad
 r\equiv(U-C)L^{-1}\pmod{2U},\quad 1\le r<2U,}{} with the exact image
\SourceMath{T^m(n)=u+2Lt,\qquad Uu=Lr+C.}{1} The comparison with the
repository's September 13 notation is the identity on the labelled word,
\(S_j=A_j\), and \(B_w=C_m\); the index substitution is \(i=j+1\). The
original coefficients, residue representative, and affine map are
identical under this comparison.

\subsection{Proposition 1. Complete first-descent fiber}

The set of original positive sources whose \textbf{first} strict
odd-return descent is at the end of exactly \(p\) is empty for
\(D\le0\). For \(D>0\) it is
\SourceMath{\{r+2Ut:t\in\mathbb Z,\ \ell\le t\le h\},}{2} where
\SourceMath{\ell=\max\!\left(0,\left\lfloor\frac{C-Dr}{2UD}\right\rfloor+1\right),
\quad
 h=\min_{\substack{j<m\\D_j>0}}
       \left\lfloor\frac{C_j-D_jr}{2UD_j}\right\rfloor.}{3} The minimum
of the empty set here is \(+\infty\). A label with \(h<\ell\) remains a
declared empty first-descent fiber; its legal cylinder and coefficients
are retained.

\textbf{Proof.} On the original cylinder,
\SourceMath{T^j(n)<n\iff D_jn>C_j.}{} For \(D_j\le0\), positivity of
\(n,C_j\) rules out strict descent. For positive \(D_j\), substituting
\(n=r+2Ut\) and using integer floors gives (3): strict descent at the
last step and non-descent at every earlier step. These are all the
conditions. The inverse source parameter is \((n-r)/(2U)\), and the
inverse target parameter is \((y-u)/(2L)\). Both are applied only on
their specified images. ∎

Distinct first-descent words have disjoint occupied fibers: the
deterministic actual orbit has one first-descent index and one exponent
list at that index. In particular, the mask for \((2,2)\) is empty,
whereas the mask for \((2)\) has \(r=1,\ell=1\); the fixed point 1 has
not been treated as a strict descent.

For every supported source in (2), retain the original edge chain
\SourceMath{b_p(n)=\sum_{j=0}^{m-1}e_{T^j(n)},\qquad
 \partial b_p(n)=v_n-v_{T^m(n)}.}{4} The relative convention is exactly
the predecessor's: only \(v_1,e_1\) are sent to zero.

\section{2. Whole exponent tails, without an overflow class}

Let \(q\) be an original word all of whose nonempty prefixes have
\(D_j\le0\). The empty word is included, with \(L=U=1,C=0,r=u=1\). Every
positive source in its cylinder has no earlier descent. Retain its chart
\SourceMath{n=r+2Ut,\quad x=T^{|q|}(n)=u+2Lt,\quad t\ge0.}{} Define
\(B\ge1\) to be the least integer satisfying the \textbf{two original
coefficient inequalities} \SourceMath{2^B r>3u+1,\qquad 2^BU\ge3L.}{5}
This integer exists. In this setting the second inequality is strict as
well: \(2^BU\) is a power of 2 and \(3L\) a positive power of 3.

\subsection{Theorem 2. Unbounded last-exponent descent family}

All sources in the prefix cylinder with actual next exponent \(a\ge B\)
have their first strict descent at that next return. Their full source
set is the single original progression
\SourceMath{n=R+2U\,2^{B-1}v,\quad v\ge0,}{6} where
\SourceMath{t_0\equiv-\frac{3u+1}{2}(3L)^{-1}\pmod{2^{B-1}},
 \quad0\le t_0<2^{B-1},\quad R=r+2Ut_0.}{7} No upper bound on \(a\) is
imposed.

The full exponent label is recovered by
\SourceMath{t=t_0+2^{B-1}v,\qquad a=\nu_2(3u+1+6Lt).}{8} Its exact
\(a\)-stratum, for each \(a\ge B\), is
\SourceMath{t\equiv\left(2^{a-1}-\frac{3u+1}{2}\right)(3L)^{-1}\pmod{2^a}.}{9}
Equation (9) gives exactly the predecessor's complete cylinder for \(q\)
followed by \(a\), with source step \(2^{A+a+1}\), target step
\(2\cdot3^{|q|+1}\), and the original affine numerator. Its
first-descent interval is the entire positive cylinder.

\textbf{Proof.} Divisibility of \(3x+1=3u+1+6Lt\) by \(2^B\) is
equivalent to (7), because \(3L\) is odd. This proves (6), with the
inverse parameter \(v=(n-R)/(2U2^{B-1})\). Specifying valuation exactly
\(a\) means congruence to \(2^a\) modulo \(2^{a+1}\); division of that
congruence by 2 gives (9). It also proves (8) and the disjoint, complete
stratification by actual exponents.

For all these sources, \SourceMath{\frac{3u+1+6Lt}{2^a}
 \le\frac{3u+1+6Lt}{2^B}<r+2Ut,}{10} where the final strict inequality
uses both inequalities in (5), including the strict constant term. Every
previous \(D_j\le0\) precludes an earlier descent. Each exact stratum
has the original word chart (1), so its intermediate values and edge
chain are retained. ∎

\subsection{A full symbolic family at the old residual frontier}

For \(q=(1,2)\), its unchanged chart is
\SourceMath{11+16t\longmapsto13+18t.}{} Equations (5)--(7) give
\(B=2,t_0=0\). Thus every \(v\ge0\) has the actual path
\SourceMath{11+32v\ \longmapsto\ 17+48v\ \longmapsto\ 13+36v
 \ \longmapsto\ \frac{10+27v}{2^{\nu_2(10+27v)}}.}{11} The full word is
\SourceMath{(1,2,2+\nu_2(10+27v)).}{} The terminal value is at most
\(10+27v<11+32v\); the first two displayed values are larger than the
source. The original three-edge chain has boundary source minus the
exact final value in (11). The residue class \(11\bmod16\) is now split
exactly: \(11\bmod32\) has this descent, while \(27\bmod32\) follows the
noncontracting child \((1,2,1)\). No unspecified probability or changed
target interval is used.

For the empty prefix, the tail is \(n=5+8v\), with \(a\ge3\). Together
with the \((2)\) cell it supplies every one-step descent source
\(n>1,n\equiv1\pmod4\). In particular,
\SourceMath{n_j=(4^j-1)/3\longmapsto1,\qquad j\ge2,}{} has actual
exponent \(2j\). For \(j\ge33\), these sources are unresolved by the
original exponent cutoff 64 but are discharged by this one actual edge,
without increasing that cutoff. The verifier executes
\(j=33,64,128,512\); the displayed identity proves the whole family.

\subsection{Finite complete next-step decomposition}

For any prefix \(q\) above, enumerate the finite list \(1\le a<B\). A
child with \(D_{qa}\le0\) remains an original prefix frontier. A child
with \(D_{qa}>0\) has its exact first-descent interval (2) and its
finite complementary non-descent box from Section 5. All \(a\ge B\)
belong to the proved tail (6). These disjoint actual-exponent cases
exhaust the full prefix cylinder. \nolinkurl{tail_descent.py} implements
this decomposition; it never adds an exponent-overflow atom.

\section{3. Attaching new paths at actual retained roots}

Let \((h,Q,F)\) be the predecessor's actual path retraction, with
\SourceMath{\partial h=I-Q,\quad hQ=0,\quad Q^2=Q,\quad F=I-h\partial.}{}
Its root \(r(n)\) is at most \(n\); \(h(v_n)\) is the original finite
path to that root. These properties hold for every whole-source cutoff
retraction in the imported code, without declaring its roots to be
convergent.

A new descent rule is applied \textbf{only at an old root} \(r\). A cell
from Section 1 or tail from Section 2 supplies an actual block
\(b:r\to d<r\). Attach the old path \(h(v_d)\), giving
\SourceMath{\widetilde b_r=b+h(v_d),\qquad
 \partial\widetilde b_r=v_r-Q(v_d).}{12} Its endpoint is an old root no
greater than \(d<r\). Iterate these root-to-root blocks. The recursion
terminates by strict decrease of positive integers. Let \(\gamma_r\) be
the resulting original edge chain and \(R(r)\) its endpoint. An unlisted
root remains itself, with zero added path.

Put
\SourceMath{K(v_n)=\gamma_{r(n)},\quad h'=h+K,\quad Q'(v_n)=v_{R(r(n))},\quad
 F'=I-h'\partial.}{13} The implementation gives an already compiled
first-descent cell priority over the identical word certified by a tail.
The original rules remain recorded. Distinct occupied first-descent
words cannot compete at a source. Tail words also end at the first
coefficient crossing, so distinct noncontracting prefixes have disjoint
tail supports.

\subsection{Theorem 3. Exact root-level extension}

These are maps on all original finite-support relative chain modules.
They satisfy
\SourceMath{\partial K=Q-Q',\quad KQ=K,\quad KQ'=0,\quad hQ'=0,}{}
\SourceMath{\partial h'=I-Q',\quad h'Q'=0,\quad Q'Q=QQ'=Q',}{}
\SourceMath{F'^2=F',\quad \partial F'=Q'\partial,\quad F'F=FF'=F',\quad FK=K.}{14}
Every closed chain is fixed literally: \(F'c=c\) for \(\partial c=0\).
The new image complex is a chain-homotopy retract of the old image
complex; its homotopy is \(K\) on the old vertex image. Original
homology and its algebraic-dual cohomology are retained by these
specific maps.

\textbf{Proof.} Telescoping (12) proves
\(\partial\gamma_r=v_r-v_{R(r)}\). Final roots have no new applicable
block and were old roots, so both homotopies vanish there. This gives
the first two lines. In particular, \(h\partial K=h(Q-Q')=0\), hence
\(FK=K\): the correction really lies in the old projected edge module.
Also \(h'\partial h'=h'(I-Q')=h'\), proving idempotence of \(F'\). The
chain equation follows from \(\partial h'=I-Q'\). Direct multiplication
using \(h'Q=K\) and \(hQ'=0\) gives both compositions with \(F\). The
difference of the old and new projections on the old image is
\((K\partial,\partial K)\), the stated chain homotopy. Substitution of
\(\partial c=0\) proves literal cycle preservation. Applying the
algebraic dual to these identities supplies the actual cochain maps and
homotopy. ∎

A nonperiodic component is not erased by an absence of cycle periods:
its degree-zero component obstruction is retained by the same homotopy
equivalence. No finite cutoff root has been promoted to such a
component.

\subsection{The 27 cylinder}

The full 37-letter word is recorded in \nolinkurl{verification.json}.
Its exact data are \SourceMath{A=59,\quad L=450283905890997363,\quad
 C=1100931843921811423,\quad D=126176846412426125.}{} They give the
entire family
\SourceMath{27+2^{60}t\longmapsto23+900567811781994726t,\qquad t\ge0.}{15}
Every proper prefix has \(D_j\le0\), and \(27D-C>0\), so (15) is a first
descent on its whole cylinder. In contrast, its original formal
reference value is
\SourceMath{F_p(3)=306472945199350439/72057594037927936>3.}{} Thus this
exact source family supplies an extension that the reference-3 stop does
not supply at this word. At cutoff 64 the old root of 27 is 27. The new
block ends at 23, whose old root is 1; (12) adds that already checked
path. The resulting 41-edge path has boundary \(v_{27}\).

\section{4. Multiplicative period morphism and finite arithmetic
exclusion}

For an actual orbit define the positive rational edge label
\SourceMath{g(n)=\frac{3n+1}{4n}=2^{a(n)-2}\frac{T(n)}n.}{16} The factor
\(T(n)/n\) is the displayed multiplicative coboundary of the actual
vertex coordinate. On a finite path it telescopes to terminal/source. On
a closed cycle its product is 1. All factors are retained; the following
inequalities are bounds on this same map, not a replacement of the
affine equations.

For a positive cycle with \(m\) odd returns, exponent sum \(A\), and
every vertex at least \(s\),
\SourceMath{3^m<2^A=\prod_{i=1}^m\left(3+\frac1{x_i}\right)
       \le\frac{(3s+1)^m}{s^m}.}{17} Let \(k=\#\{i:a_i=1\}\). Since
\SourceMath{A=2m-k+E,\qquad E=\sum_{a_i\ge3}(a_i-2)\ge0,}{} (17) implies
the exact integer inequality \SourceMath{(4s)^m\le2^k(3s+1)^m.}{18} The
projection to \((m,A,k,s)\) retains its fiber: all original labelled
words with these counts whose unique primitive \(x_i=C(p_i)/D\) is
integral, satisfies \(x_i\ge s\), and satisfies the declared minimum
bound. The matrix forcing and word order in this fiber are not declared
solvable merely because a window (17) is allowed.
\nolinkurl{integral_blocks.md} gives the integral comparison maps for
this remaining fiber.

\subsection{Period-driven feedback actually executed}

The atlas initially verifies an actual first descent for every odd
source \(3\le n\le4095\). Strong induction on the source, beginning at
1, proves that all these sources reach 1. Consequently every nontrivial
cycle has minimum at least 4097.

For each current lower bound \(s\), select the least period for which
(17) has any integer exponent \(A\). The least possible exponent is the
bit length of \(3^m\). For that pair, compute the largest integer \(u\)
satisfying \SourceMath{2^Au^m\le(3u+1)^m.}{19} The ratio \((3+1/u)^m\)
strictly decreases with positive \(u\), so a successful comparison at
\(u\) and failed comparison at \(u+1\) prove the exact ceiling. The
original minimum of this entire period window lies in a finite interval.
The executed continuation proves first descent throughout that interval,
then recomputes the next window.

{\def\LTcaptype{none} % do not increment counter
\begin{longtable}[]{@{}
  >{\raggedleft\arraybackslash}p{(\linewidth - 6\tabcolsep) * \real{0.2500}}
  >{\raggedleft\arraybackslash}p{(\linewidth - 6\tabcolsep) * \real{0.2500}}
  >{\raggedleft\arraybackslash}p{(\linewidth - 6\tabcolsep) * \real{0.2500}}
  >{\raggedleft\arraybackslash}p{(\linewidth - 6\tabcolsep) * \real{0.2500}}@{}}
\toprule\noalign{}
\begin{minipage}[b]{\linewidth}\raggedleft
Current minimum lower bound
\end{minipage} & \begin{minipage}[b]{\linewidth}\raggedleft
First window \((m,A)\)
\end{minipage} & \begin{minipage}[b]{\linewidth}\raggedleft
Exact minimum ceiling
\end{minipage} & \begin{minipage}[b]{\linewidth}\raggedleft
New verified odd-source bound
\end{minipage} \\
\midrule\noalign{}
\endhead
\bottomrule\noalign{}
\endlastfoot
4097 & (147,233) & 6724 & 6723 \\
6725 & (200,317) & 12824 & 12823 \\
12825 & (253,401) & 27113 & 27113 \\
27115 & (306,485) & 99780 & 99779 \\
\end{longtable}
}

Every new seed is iterated only to its actual first lower odd value.
Independent repeated-division replay verifies the entire final interval,
and strong induction supplies its boundary paths to 1. These seeds
produce \textbf{5,175 occupied infinite first-descent cells}; their
symbolic lift statements follow from Proposition 1, not from sampling.
The construction additionally compiles \textbf{4,011 unbounded
last-exponent tails} from the admitted prefix words.

\subsection{Theorem 4. Certificate excluding all cycles with at most 402
exponent-1 steps}

There is no nontrivial positive integer Collatz cycle whose primitive
odd-return word contains at most 402 occurrences of exponent 1.

\textbf{Proof.} The actual seed certificates above establish
\(s=99781\). The following finite integer comparisons are executed, with
an independent least-power-of-two loop in the verifier:
\SourceMath{2^{\operatorname{bitlength}(3^m)}s^m>(3s+1)^m
 \quad(1\le m\le970).}{20} They exclude every integer exponent in (17)
at those lengths. Separately,
\SourceMath{(4s)^{969}>2^{402}(3s+1)^{969}.}{21} Since \(4s>3s+1>0\),
the same strict comparison holds at every larger length with any
\(k\le402\). Thus (18) forces \(m\le968\), whereas (20) excludes all
those periods. The finite premises are part of the executable
certificate; neither a sampled frequency nor an unproved analytic
estimate occurs. ∎

This is a self-contained improvement of the present workbench's previous
zero-or-one-exponent-1 exclusion, \textbf{not a claimed advance over
published cycle bounds}. The executed feedback used a seed budget of
100,000; that budget restricts the computation, not the all-length scope
of Theorem 4.

The first unexcluded window at the new minimum is
\SourceMath{(m,A,k_{\min})=(971,1539,403),\qquad
 99781\le\min x_i\le330749.}{22} For \(k=403\), (18) forces \(m\le971\),
and (20) forces \(m=971\). The only exponent in (17) at this length is
1539, as the next power of 2 fails the upper bound. Then \(E=0\): the
original word consists of 403 ones and 568 twos. This window is retained
as an unresolved arithmetic fiber, not as a found cycle. Its integral
compression is in \nolinkurl{integral_blocks.md}.

The precise clock comparison with conventional shortened-map
local-minimum counts is also explicit. An odd-return exponent \(a\)
expands to one odd shortened step followed by \(a-1\) even shortened
steps. Expansion sends total shortened length to \(A\) and odd-step
count to \(m\); grouping between successive odd states is its inverse.
Each \(a\ge2\) contributes exactly one even-run termination and hence
one local minimum in a nontrivial shortened cycle. Thus that count is
\(\mu=m-k\). The original ordered word remains available. This supplies
the map between these counts rather than interchanging them with the
\(m\)-cycle terminology in {[}R2{]}.

\section{5. First-crossing exception boxes and the retained infinite
source}

At a first \textbf{completed odd-return coefficient crossing}, \(D>0\)
and every proper prefix has \(D_j\le0\). Its exact positive sources with
no descent even at the terminal step are \SourceMath{n=r+2Ut,\qquad
 \max\!\left(0,\left\lceil\frac{s-r}{2U}\right\rceil\right)
 \le t\le\left\lfloor\frac{C-Dr}{2UD}\right\rfloor.}{23} This is the
complement of the descent interval within that original cylinder above
\(s\). It is a finite \textbf{source} box, not an absent label. The
original affine offset is essential to its upper endpoint.

There is a general bound on its occupancy:
\SourceMath{3C\le mL,\qquad \#\text{box}\le\left\lceil\frac{m}{6D}\right\rceil.}{24}
Indeed each summand of \(C\) is at most \(3^{m-1}\), using
\(2^{A_{i-1}}\le3^{i-1}\). The possible sources lie in \((0,C/D]\), are
spaced by \(2U\), and \(C/(2UD)\le mL/(6UD)<m/(6D)\). This proves (24).
It does not assert the box to be empty.

Every first-crossing word with total exponent sum at most 18 was
checked: 4,082 words, 4,081 empty boxes at minimum 1, and one occupied
box containing the actual fixed source 1. This is a bounded scan. The
decomposition in Section 2 applies at arbitrary finite prefixes, without
asserting that their infinite-depth residual source is empty.

An actual non-base component minimum has no strict descent. It must
avoid every certified cell and tail. A finite coefficient crossing
places it in exactly (23); absence of any coefficient crossing retains
all the original noncontracting prefix cylinders. In both cases ordinary
positive-integer support still requires the predecessor's canonical
residue tower to stabilize at that same integer. No reference-measure
zero is used to remove a vertex. The root splicing in Section 3
preserves the original component obstruction in either case.

The distinction between coefficient contraction and actual descent is
connected to the explicit remainder in {[}R1{]} by the clock expansion
above: at a completed return, its coefficient is \(L/U\) and its
remainder is \(C/U\). Their expression \((L/U)n+C/U\) is exactly (1). No
coefficient-stopping-time conjecture is assumed here.

\section{6. Reproduction and inspection scope}

\nolinkurl{verify.py} regenerates the atlas and feedback, checks source
lifts by independent integer repeated division, checks the exact
root-level homotopies, exercises the unbounded tail strata, checks
integral block comparisons, and verifies the finite all-length exclusion
certificate. It also retains zero masks and rejects malformed data.
\nolinkurl{replay.py} checks source hashes, replays the entire unchanged
import, and requires the new ordinary/optimized outputs to be
byte-identical to \nolinkurl{verification.json}.

The current-source reads include the workbench method chapter, issue
\#3, PR \#2 and its discussion, and the cycle note at the pinned commit.
The uploaded publication archive supplied executable local bytes; its
note Git blob was compared with the authenticated source, and the
unchanged 21-file manifest and six replay outputs passed before
continuation. No unpushed local Codex state was accessed.

\textbf{{[}P{]}}
\nolinkurl{../cycle_relative_cohomology_20260914/note.md}, Sections
1--7, especially the actual relative graph and cutoff retractions;
\nolinkurl{../stopped_affine_transport_20260914/} is retained through
that original dependency.

\textbf{{[}R1{]}} O. Rozier and C. Terracol, \emph{Paradoxical behavior
in Collatz sequences}, arXiv:2502.00948v2 (13 February 2025),
Introduction, Definitions 1.1--1.2 and their displayed affine
expression. Readable primary text:
\nolinkurl{https://arxiv.org/html/2502.00948v2}. Cited for the existing
coefficient/remainder question, not as a proof premise or an adoption of
a conjecture.

\textbf{{[}R2{]}} C. Hercher, \emph{There are no Collatz m-Cycles with m
\textless= 91}, Journal of Integer Sequences 26 (2023), Article 23.3.5.
Publisher page:
\nolinkurl{https://cs.uwaterloo.ca/journals/JIS/VOL26/Hercher/hercher5.html},
which also links a corrigendum dated 14 June 2026. The page and
terminology were inspected; no detailed numerical theorem or external
verification bound from it is a premise. Its TeX retrieval was not
readable through the web tool. No priority or cycle-record claim is
made.
